\chapter{Layer L3: deterministic parametrix}\label{ch:detparametrix}

The deterministic content of the paper's \S3.1--\S3.2 and all of \S4:
profile kernels, the renormalization induction, the constants
\(\mathcal{C}_{2q}\), and the deterministic halves of Proposition~3.5.

\begin{definition}[Deterministic profile kernels; paper (3.1)--(3.2)]\label{D-Kdet}
  \uses{D-pair,D-int,I-green}
  \lean{Anderson4D.detIntegrand, Anderson4D.assemble,
    Anderson4D.SmoothCutoff.etaEpsT4}\leanok
  For each \((m,\kappa)\), the deterministic factor of the chaos
  expansion: products of \(G\)-factors along the chain and
  \(\eta_\varepsilon\)-factors along the pairs, with the single-index
  variables as free arguments; generalized inputs
  \([G_0,\dots,G_m]\).
\end{definition}

\begin{definition}[Renormalization induction; paper Def.~3.1]\label{D-RI}
  \uses{D-int,D-Kdet}
  \lean{Anderson4D.detRIkernel, Anderson4D.detRIrecursiveFull,
    Anderson4D.detRIrecursiveProfile}\leanok
  \(\mathcal{RI}\) at the deterministic level: repeatedly remove the
  smallest-leftmost fully paired subinterval, via \(J\) and
  \(\widetilde G\) ((3.3)--(3.5)).
\end{definition}

\begin{proposition}[Closed formula; paper Prop.~3.2]\label{P-3.2}
  \uses{D-RI}
  \lean{Anderson4D.raw_sub_counterterm_eq_difference,
    Anderson4D.detRIkernel_eq_prod,
    Anderson4D.detRIrecursiveFull_eq_detRIfull,
    Anderson4D.detRIrecursiveProfile_eq_detRIprofile}\leanok
  The explicit formula (3.6): difference factors
  \(G(x_{r_i} - x_{r_i+1}) - G(x_{\ell_i} - x_{r_i+1})\) at the removed
  endpoints.
\end{proposition}

\begin{definition}[Renormalization constants; paper (3.8)--(3.11)]\label{D-C2q}
  \uses{D-int,D-RI}
  \lean{Anderson4D.detJ, Anderson4D.renormC2q,
    Anderson4D.renormCEps}\leanok
  \(\mathcal{J}_{2q,\sigma}\), \(\mathcal{C}_{2q,\sigma}\),
  \(\mathcal{C}_{2q}\), and
  \(\mathcal{C}_\varepsilon = \sum_{q\le A}\mathcal{C}_{2q}\); the
  primitive-with-insertions structure of full pairings.
\end{definition}

\begin{proposition}[Closed formula for \(\mathcal{J}\); paper Prop.~3.3]\label{P-3.3}
  \uses{D-C2q,P-3.2}
  \lean{Anderson4D.renormJ_eq_prod}\leanok
  The explicit formula (3.12).
\end{proposition}

\begin{proposition}[Bound on the constants; paper (3.22)]\label{P-3.5a}
  \uses{D-C2q,P-4.1,R-322}
  \lean{Anderson4D.R322AnalyticResidualPrefix.exists_r322_renormC2q_bound}\leanok
  \(\lvert\mathcal{C}_{2q}\rvert \le \varepsilon^{-2}
  \lvert\log\varepsilon\rvert^{-1}(C\lambda)^{2q}\) for \(1 \le q \le A\).
\end{proposition}

\begin{lemma}[Reduction machine for (3.22); paper \S4.1]\label{R-322}
  \uses{P-4.1,I-green,D-int,D-C2q}
  \lean{Anderson4D.R322AnalyticResidualPrefix.exists_r322RenormFiberReductionOutputAE,
    Anderson4D.R322AnalyticResidualPrefix.exists_terminalContext_endpointFiberDetJSum_eq}\leanok
  Dyck-word bookkeeping of nested fully paired subintervals; iterative
  reduction integrating out innermost intervals; Taylor expansion with
  \(\mathcal{E}\)-symmetry killing the linear term (4.9); the three-case
  integral bounds (4.10)--(4.12).
\end{lemma}

\begin{proposition}[Deterministic second-moment bound; paper \S4.2]\label{P-3.5b-det}
  \uses{P-4.1,R-324}
  \lean{Anderson4D.deterministic_second_moment_bound}\leanok
  The deterministic pairing-sum bound underlying (3.24), including the
  decay factors
  \(\langle\alpha\rangle^{-4}\langle\beta\rangle^{-4}
  \langle\varepsilon^2(\alpha+\beta)\rangle^{-8}\).
\end{proposition}

\begin{lemma}[Reduction machine for (3.24); paper \S4.2]\label{R-324}
  \uses{P-4.1,R-322,D-RI}
  \lean{Anderson4D.R324WithinHalfResidualPrefix.exists_r324PaperResidualEndpointWeightedMajorantBound,
    Anderson4D.exists_r324PaperFullEndpointZeroShiftWeightedMajorantBound,
    Anderson4D.SmoothCutoff.exists_r324PaperHighWholeSeriesWeightedMajorantBound}\leanok
  The same reduction machine plus oscillation in \((x,y,z,w)\) for the
  decay factors ((4.16)--(4.20)).  The formal proof follows the paper's
  two endpoint cases and retains the truncation
  \(A = \lfloor|\log\varepsilon|\rfloor\) until the final numerical
  assembly.

  In Step~4(B), the paper obtains the total-frequency gain by finding a
  first large/high slot among the at most \(A\) composition factors.  The
  Lean proof keeps the complete signed open series through the removals,
  translates the common surviving left half, extracts its Fourier
  coefficient, selects a nonzero coordinate of \(\alpha+\beta\), and
  integrates by parts eight times.  The only norm is taken after this
  whole-family Fourier extraction.  This is an alternative Lean proof of
  the same required
  \(\langle\varepsilon^2(\alpha+\beta)\rangle^{-8}\) decay conclusion,
  not the paper's literal first-large/high-slot proof.

  Declaration-by-declaration links are recorded in
  \href{https://github.com/sairmath/anderson4d_lean/blob/main/docs/PAPER_TO_LEAN.md#step-4b}
  {\texttt{PAPER\_TO\_LEAN.md}, Step~4(B)}.
\end{lemma}
